\documentclass[11pt,a4paper]{article}

\usepackage{amsmath,amssymb,amsthm}
\usepackage{hyperref}
\usepackage[numbers,sort&compress]{natbib}
\usepackage[margin=1in]{geometry}
\usepackage{xcolor}

\newtheorem{theorem}{Theorem}[section]
\newtheorem{conjecture}[theorem]{Conjecture}
\newtheorem{corollary}[theorem]{Corollary}
\newtheorem{lemma}[theorem]{Lemma}
\newtheorem{proposition}[theorem]{Proposition}
\newtheorem{definition}[theorem]{Definition}
\theoremstyle{remark}
\newtheorem{remark}[theorem]{Remark}

\newcommand{\PT}{\mathbb{PT}}
\newcommand{\CP}{\mathbb{CP}}
\newcommand{\C}{\mathbb{C}}
\newcommand{\Z}{\mathbb{Z}}
\newcommand{\R}{\mathbb{R}}
\newcommand{\cO}{\mathcal{O}}
\newcommand{\cC}{\mathcal{C}}
\newcommand{\cN}{\mathcal{N}}
\newcommand{\cU}{\mathcal{U}}
\newcommand{\PP}{\mathcal{P}}
\newcommand{\barPP}{\overline{\mathcal{P}}}


\title{Twistor Theory from the Obstruction Ladder:\\
The Googly Problem as an $H^2$ Obstruction\\
{\large (Paper VII)}}


\author{Zhou-Li Chen \\ co-nlang Research}
\date{May 2026}

\begin{document}

\maketitle

\begin{abstract}
  We identify the cohomological mechanism underlying Penrose's
  googly problem---the inability of standard twistor theory to
  describe both chiralities of massless fields---as an $H^2$
  obstruction in the sense of the cohomological obstruction ladder
  (Papers~\cite{PaperI,PaperIII,PaperIV,PaperVI}), and construct the
  algebraic framework for its resolution. The central construction is a functor
  $\Phi: \mathcal{N}(\mathcal{C}(\mathcal{P}_2)) \to \CP^3$
  from the \v{C}ech nerve of the Pauli MASA poset to twistor space.
  We show that the non-trivial central extension class
  $[f] \in H^2(\bar{\mathcal{P}}_2, \Z/2)$ (Paper~\cite{PaperIII}) maps under $\Phi^*$
  to $c_1(\mathcal{O}(1)) \bmod 2 \in H^2(\CP^3, \Z/2)$, the mod-2 reduction
  of the first Chern class of the twistor spinor bundle.
  This identifies the algebraic obstruction to chirality unification
  as the same $\Z/2$ class that distinguishes $SO(3)$ from $SU(2)$.

  A central difficulty is that the standard spectral sequence for the
  Penrose transform has an empty source for the $d_2$ differential
  when computed with any single-space coefficient sheaf---both
  $\underline{\Z/2}$ (vanishes by simple-connectivity) and $\cO^*$
  (integer grading mismatch) fail to populate $E_2^{0,1}$.
  We argue that this coefficient paradox is the algebraic reason
  the googly problem has persisted for nearly six decades, and that its
  resolution requires the cross-space functor $\Phi$, which identifies
  the normal bundle of the twistor curve with the 2-qubit Hilbert space
  and transfers the PM obstruction to a cup product pairing on the
  mod-2 reduced deformation lattice.
  The chirality-mixing map $d_2^{\text{eff}}$ is thus not an
  intra-spatial spectral sequence differential but an intersection
  pairing driven by quantum contextuality.
\end{abstract}

%------------------------------------------------------------------
\section{Introduction}
\label{sec:intro}
%------------------------------------------------------------------

Penrose's twistor theory \cite{Penrose1967} provides a remarkable
encoding of massless fields as sheaf cohomology classes on
$\PT \cong \CP^3$:
\[
  H^1(\PT^+, \cO(-n-2)) \;\cong\;
    \text{massless fields of helicity } n/2.
\]

This isomorphism succeeds spectacularly for self-dual
(positive helicity) fields. For anti-self-dual (negative helicity)
fields---the \emph{googly} sector in Penrose's cricket terminology
\cite{PenroseRindler}---no such unified description exists on a
single twistor space. For nearly six decades, this remains the central
open problem of twistor theory~\cite{Penrose2004}.

Prior approaches can be grouped into three strategies. The
\textit{ambitwistor} programme \newline \cite{MasonSkinner2014,CachazoHeYuan2014}
achieves chirality unification by working on the product space
$\PT \times \PT^*$, effectively comparing both helicity sectors
simultaneously via a worldsheet model. The \textit{palatial twistor}
programme~\cite{Penrose2004,Penrose2016} seeks a larger twistor space with a
non-linear structure that accommodates both signs of helicity,
but lacks an explicit cohomological mechanism. The \textit{twistor
action} framework~\cite{Adamo2017} computes amplitudes directly from
holomorphic Chern-Simons theory on $\PT^+$ but does not resolve the
asymmetric treatment of helicities at the cohomological level.
The \textit{topos-theoretic} approach to quantum
contextuality~\cite{IshamButterfield1999} provides the sheaf-theoretic
language that makes the MASA poset construction natural, but has not
previously been connected to the googly problem.

The present paper pursues a fourth strategy: the googly problem is
not a failure of $\PT$ as a space, but a failure to recognise it
as an $H^2$ obstruction of the same type that separates $SO(3)$ from
$SU(2)$.

The cohomological obstruction ladder~\cite{PaperI,PaperIII,PaperIV,PaperVI} develops a
spectral sequence framework for quantum contextuality, built on
the \v{C}ech nerve of the maximal abelian subalgebra (MASA) poset
$\cC(A)$. This framework classifies obstructions to classical
descriptions by cohomological degree:
\begin{itemize}
  \item $H^1$: geometric phases ($U(1)$ holonomy) --- Paper I~\cite{PaperI}
  \item $H^2$: Kochen-Specker~\cite{KochenSpecker} contextuality (central extensions by
    $\Z/2$) --- Paper III~\cite{PaperIII}
  \item $H^3$: Borromean contextuality (Dixmier-Douady gerbes) ---
    Paper IV~\cite{PaperIV}
\end{itemize}

The present paper demonstrates that these two independently
developed frameworks---twistor theory and the obstruction ladder---
are structurally united. Specifically:
\begin{enumerate}
  \item The googly problem is an $H^2$ obstruction: $H^1$ can encode
    phase but not chirality, which requires the $\Z/2$ data carried
    by $H^2$.
  \item The mechanism is identical to the $SO(3) \to SU(2)$ central
    extension: in both cases, a non-trivial class
    $[f] \in H^2(\cdot, \Z/2)$ unlocks a sector invisible at the
    $H^1$ level.
  \item We construct the critical mapping $\Phi$ from the Pauli MASA
    \v{C}ech nerve to $\CP^3$, via an augmented cover $\mathcal{W}$
    whose patches are explicitly non-contractible ($\tilde{U}_i \simeq S^2$).
    The identification $\Phi^*([f]) = c_1(\cO(1)) \bmod 2$ is
    Proposition~\ref{conjecture:global_embedding}, requiring explicit cochain-level verification.
  \item The classical spectral sequence collapses because its $d_2$
    source is empty: both $\underline{\Z/2}$ and $\cO^*$ fail as
    coefficient sheaves for $E_2^{0,1}$. We identify this coefficient
    paradox as the algebraic core of the decades-old googly problem,
    and resolve it using the cross-space functor $\Phi$ via an
    integral lattice construction: the normal bundle's integer sections
    $Ext^1_{\Z} \cong \Z^4$ are mod-2 reduced to $(\Z/2)^4$, on which
    $\Phi^*([f])$ acts by cup product. The $d_2^{\text{eff}}$ map is
    thus an intersection pairing on the mod-2 reduced deformation
    lattice, not an intra-spatial differential operator.
  \end{enumerate}

\subsection*{Logical dependency}

The argument proceeds in the following order:
\[
  \text{Construct } \Phi \;\to\;
  \text{Pull back MASA } d_2 \neq 0 \;\to\;
  \text{Both chiralities appear at } E_3.
\]

Paper III~\cite{PaperIII}'s $d_2 \neq 0$ is computed on the MASA nerve, not on $\PT$.
The googly resolution requires $\Phi$ to transfer this non-vanishing
class. The $SU(2)/\CP^1$ computation---developed in \S\S 3--5---is
precisely this transfer test.

%------------------------------------------------------------------
\section{The googly problem as an $H^2$ obstruction}
\label{sec:googly-h2}
%------------------------------------------------------------------

\subsection{Chirality and cohomological degree}

The Penrose transform occupies $H^1$---the level of $U(1)$ phase.
This is the cohomological degree that encodes Berry-like geometric
phases (Paper I~\cite{PaperI}). But chirality is a $\Z/2$ distinction:
it distinguishes orientation, not accumulated phase.

In the obstruction ladder hierarchy:
\begin{itemize}
  \item $H^1$ classes are $U(1)$-valued. They classify line bundles.
    They carry holonomy information but no orientation data.
  \item $H^2$ classes can be $\Z/2$-valued. They classify central
    extensions. They carry $\pm 1$ sign data, which is precisely
    the distinction between left and right chirality.
\end{itemize}

The googly problem is therefore not a failure of twistor theory per
se, but a consequence of operating at $H^1$. Anti-self-dual fields
are $H^2$ objects attempting to live in an $H^1$ framework.

\subsection{The $SO(3) \to SU(2)$ parallel}

The most precise model for this situation is the central extension
\[
  1 \to \Z/2 \to SU(2) \to SO(3) \to 1,
\]
classified by $[f_{SO}] \in H^2(SO(3), \Z/2) = \Z/2$.

\begin{itemize}
  \item $SO(3)$ has only integer-spin representations. A physicist
    working solely with $SO(3)$ would conclude that half-integer
    spin states do not exist.
  \item $SU(2)$ is the non-trivial central extension. It unlocks
    half-integer spin representations.
  \item Electrons are physically real---and they are invisible from
    the $SO(3)$ perspective.
\end{itemize}

We propose an exact parallel in twistor theory:
\begin{itemize}
  \item $\PT \cong \CP^3$ has only self-dual fields. A physicist
    working solely with $\PT$ concludes that anti-self-dual fields
    do not exist.
  \item $\widetilde{\PT}$---a central extension of $\PT$ by
    $[e] \in H^2(\PT, \Z/2)$---unlocks both chiralities.
  \item Anti-self-dual fields are physically real---and they are
    invisible from the $\PT$ perspective.
\end{itemize}

Paper III~\cite{PaperIII} computed precisely such a central extension:
the Peres-Mermin~\cite{PeresMermin} square yields $[f] \in H^2(G_{\text{PM}}, \Z/2)$
with the 4-cycle accumulating $-\mathbf{I}$. The central claim of
this paper is that $[f]$ maps to $c_1(\cO(1)) \bmod 2$ under the
$\Phi$ functor (\S\ref{sec:phi-construction}).

%------------------------------------------------------------------
\section{Pauli MASA cover of $\CP^1$}
\label{sec:masa-cp1}
%------------------------------------------------------------------

\subsection{Three contextual MASA patches from the Peres-Mermin~\cite{PeresMermin} square}

Instead of choosing naive axis-aligned 1-qubit subalgebras---which
automated cohomology computation~(\cite{ScriptConj42}) confirms
will split under global mod-2 reduction---we extract three maximal
abelian subalgebras directly from the columns of the
Peres-Mermin~\cite{PeresMermin} magic square. Each column defines a
distinct commuting context of
operators in the 2-qubit Pauli algebra $\PP_2$:
\begin{align*}
  \mathfrak{h}_1 &= \langle X \otimes I,\; I \otimes Z,\; X \otimes Z \rangle,\\
  \mathfrak{h}_2 &= \langle I \otimes X,\; Z \otimes I,\; Z \otimes X \rangle,\\
  \mathfrak{h}_3 &= \langle X \otimes X,\; Z \otimes Z,\; Y \otimes Y \rangle.
\end{align*}
Each $\mathfrak{h}_i$ forms a maximal commutative subalgebra whose
simultaneous eigenbases define mutually unbiased bases of $\C^4$.
While any individual column consists of strictly commuting elements,
their global cross-contextual entanglement constitutes the exact
topological backbone of the Kochen-Specker~\cite{KochenSpecker} contradiction.

\subsection{Geometric realization on $\CP^1$}

Let $\Sigma_0 \cong \CP^1$ be the twistor curve corresponding to
the spacetime origin $x = 0$, parameterized by homogeneous coordinates
$[Z^0 : Z^1]$. We establish a discrete open cover
$\cU = \{U_1, U_2, U_3\}$ on $\Sigma_0$, mapping each column MASA
$\mathfrak{h}_i$ to a distinguished spatial sector on the Riemann
sphere:
\begin{align*}
  U_1 &= \{[Z^0:Z^1] \mid |Z^1/Z^0| < 1 + \varepsilon\}
    \quad (\text{neighborhood of } 0 \leftrightarrow \mathfrak{h}_1),\\
  U_2 &= \{[Z^0:Z^1] \mid |Z^0/Z^1| < 1 + \varepsilon\}
    \quad (\text{neighborhood of } \infty \leftrightarrow \mathfrak{h}_2),\\
  U_3 &= \CP^1 \setminus \{\text{small disks around } 0, \infty\}
    \quad (\text{equatorial annulus} \leftrightarrow \mathfrak{h}_3).
\end{align*}
The \v{C}ech nerve $\cN(\cU) = \cN(\{U_1, U_2, U_3\})$ generated
by this assignment maps the overlapping intersections directly onto
the relational transitions of the Peres-Mermin~\cite{PeresMermin} configuration:
\begin{itemize}
  \item 0-simplices: $\{U_1, U_2, U_3\}$ (the three columns)
  \item 1-simplices: $\{U_{12}, U_{23}, U_{31}\}$ (pairwise overlaps)
  \item 2-simplex: $U_{123} = U_1 \cap U_2 \cap U_3$ (triple intersection)
\end{itemize}

\subsection{\v{C}ech 1-cocycle from L-S contraction}

We define the local transition functions $g_{ij}: U_{ij} \to U(1)$
via the holonomy of the Lohmiller-Slotine (L-S) contraction metrics
as the system undergoes a continuous gauge transition between
different column contexts:
\[
  g_{12} = e^{i\theta_{12}},\quad
  g_{23} = e^{i\theta_{23}},\quad
  g_{31} = e^{i\theta_{31}},
\]
where $\theta_{ij}$ represents the geometric phases accumulated by
parallel transport across the overlapping patches on the nerve.
The \v{C}ech differential $\delta: C^1 \to C^2$ evaluates their
closure on the triple intersection:
\[
  (\delta g)_{123}
    = g_{12} \cdot g_{23} \cdot g_{31}
    = \exp\!\big(i(\theta_{12} + \theta_{23} + \theta_{31})\big).
\]

These $U(1)$-valued transition data must be promoted to
$SO(3)$-valued clutching functions for the transgression argument
of \S\ref{sec:conj42-transgression}.  The promotion proceeds in two
steps:
\begin{enumerate}
  \item Embed each $g_{ij} \in U(1)$ into $SU(2)$ as a diagonal
    matrix $\tilde g_{ij} = \operatorname{diag}(g_{ij}, g_{ij}^{-1})$.
    The PM column-3 product
    $(X \otimes X)(Z \otimes Z)(Y \otimes Y) = -\mathbf{I} \in SU(2)$
    corresponds to $\tilde g_{12} \tilde g_{23} \tilde g_{31}$.
  \item Project via the adjoint representation
    $\operatorname{Ad}: SU(2) \to SO(3)$, whose kernel is $\Z/2
    = \{\pm\mathbf{I}\}$.  The image
    $\operatorname{Ad}(\tilde g_{ij})$ is an $SO(3)$-valued
    transition function on $U_{ij}$, and the triple product
    $-\mathbf{I} \in SU(2)$ maps to the non-trivial element of
    $\pi_1(SO(3)) \cong \Z/2$, generating the clutching class
    used in \S\ref{sec:conj42-transgression}.
\end{enumerate}
The projection is well-defined because the PM column-3 configuration
is symmetric under the $\Z/2$ kernel: $-\mathbf{I}$ is
distinguished by its non-trivial class in $\pi_1(SO(3))$, not by
its lift to $SU(2)$.

\subsection{The configurational $H^2$ obstruction}

A critical subtlety unveiled by automated group cohomology
verification~(\cite{ScriptConj42}) is that the non-abelian nature
of the full Pauli group is insufficient to secure a stable
obstruction: the global mod-2 reduction of
$H^2(\barPP_2/Z, \Z/2)$ splits completely.
The un-absorbable topological anomaly is strictly
\textit{configurational}, driven by the explicit row-column
constraint of the Peres-Mermin~\cite{PeresMermin} grid:
\[
\begin{array}{ccc}
  X \otimes I & I \otimes X & X \otimes X \\
  I \otimes Z & Z \otimes I & Z \otimes Z \\
  X \otimes Z & Z \otimes X & Y \otimes Y
\end{array}
\]

All row products close symmetrically to $+\mathbf{I}$, but the
columns host a structural contradiction. Evaluating the continuous
transitions through the first two columns yields trivial accumulated
phases ($+\mathbf{I}$), whereas the third column accumulates the
central sign defect:
\[
  (X \otimes X)(Z \otimes Z)(Y \otimes Y) = -\mathbf{I}.
\]

Because the third column is mapped directly onto the equatorial
intersection $U_{123}$ that governs the triple overlap on the
$\CP^1$ nerve, the L-S transition functions are structurally
prevented from closing to unity:
\[
  (\delta g)_{123} = \omega = -1.
\]

This $\omega \neq 1$ is a rigid \v{C}ech 2-cocycle representing a
non-trivial restricted class
$[f] \in H^2(G_{\text{PM}}, \Z/2)$.
It bridges quantum contextuality to twistor geometry not as a
loose algebraic property of the operators, but as an inescapable
topological defect of the measurement configuration itself.

\subsection{Why $\CP^1$ cannot absorb this class}

The obstruction vanishes in $H^2(\CP^1, \cO^*)$ because
$H^2(\CP^1, \Z) = 0$. The exponential sequence
\[
  0 \to \Z \to \cO \to \cO^* \to 0
\]
gives $H^2(\CP^1, \cO^*) \cong H^3(\CP^1, \Z) = 0$.
There is no room for a non-trivial 2-cocycle in the holomorphic
geometry of a single Riemann sphere.

A subtle but essential point: $H^2(\CP^1, \Z/2) = \Z/2$ by the
universal coefficient theorem, so $\CP^1$ \emph{can} accommodate this
class \emph{topologically}. The obstruction is \emph{holomorphic}:
$H^2(\CP^1, \cO^*) = 0$ means no \emph{holomorphic} line bundle or
gerbe can carry the class. The googly problem is thus a specifically
\emph{holomorphic} obstruction: anti-self-dual fields require a
holomorphic $H^2$ class, which $\CP^1$ alone cannot supply.

This is the geometric content of the googly problem:
the $\CP^1$ at a spacetime point is a holomorphically shallow
space that cannot host the $H^2$ obstruction required for
chirality unification. The next section constructs the lift to
$\CP^3$, where $H^2(\CP^3, \cO^*) \ni c_1(\cO(1)) \neq 0$
provides the necessary holomorphic cohomological capacity,
via Serre's GAGA theorem \cite{SerreGAGA}.

%------------------------------------------------------------------
\section[Global Embedding and Cover Non-Triviality]
  {Global embedding and cover non-triviality}
\label{sec:phi-construction}
%------------------------------------------------------------------

{\small\emph{Note: The construction in this section is conjectural.}
While the MASA cover of $\CP^1$ (\S\ref{sec:masa-cp1}) is explicit
and the triple overlap computation is verified, the lift to $\CP^3$
(via the augmented cover $\mathcal{W}$) and the claim
$\Phi^*([f]) = c_1(\cO(1)) \bmod 2$ require verification by direct
spectral sequence computation. The GAGA comparison (\S\ref{sec:gaga})
is mathematically rigorous but depends on the lift existing.}

To establish the comparison map $\Phi^*$ rigorously, we translate
the algebraic contextuality class $[f] \in H^2(\barPP_2, \Z/2)$
into a concrete \v{C}ech 2-cocycle on the projective twistor space
$\PT = \CP^3$.

Let $\Sigma_0 \subset \CP^3$ be the holomorphic $\CP^1$ curve
corresponding to the spacetime origin $x = 0$, defined by the
vanishing of the transverse homogeneous coordinates:
\begin{equation}
  \Sigma_0 = \bigl\{ [Z^0 : Z^1 : Z^2 : Z^3] \in \CP^3
    \;\big|\; Z^2 = 0,\; Z^3 = 0 \bigr\}.
\end{equation}
On $\Sigma_0$ the standard inhomogeneous coordinate is
$\xi = Z^1/Z^0$. Following Paper III~\cite{PaperIII}, we cover $\Sigma_0$ with
three overlapping MASA patches
$\mathcal{U} = \{U_1, U_2, U_3\}$ as in \S\ref{sec:masa-cp1}.

We construct the global cover of $\CP^3$ by performing a cylindrical
extension of the MASA patches along the transverse directions,
regularized by the standard remaining charts of $\CP^3$.
Define the augmented \v{C}ech cover
$\mathcal{W} = \{\tilde{U}_1, \tilde{U}_2, \tilde{U}_3, V_2, V_3\}$,
where for $i \in \{1,2,3\}$,
\begin{equation}
  \tilde{U}_i = \bigl\{ [Z^\alpha] \in V_0 \cup V_1
    \;\big|\; [Z^0:Z^1:0:0] \in U_i \bigr\},
\end{equation}
and $V_2 = \{Z^2 \neq 0\}$, $V_3 = \{Z^3 \neq 0\}$.

\begin{remark}[Non-good cover]
\label{rem:non-good}
The open patches $\tilde{U}_i$ are not contractible. Each removes a
projective line $\ell_i \cong \CP^1$ from $\CP^3$, yielding the
homotopy type:
\[
\tilde{U}_i \simeq \CP^3 \setminus \ell_i \simeq \CP^1 \cong S^2.
\]
Since $H^2(S^2, \Z) = \Z \neq 0$, the standard isomorphism between
\v{C}ech cohomology $\check{H}^2(\mathcal{W}, \Z/2)$ and global
sheaf cohomology $H^2(\CP^3, \Z/2)$ does not follow from the
classical Leray theorem for good covers. A generalized Leray theorem
for non-contractible patches, given by the pre-sheaf spectral
sequence (Godement~\cite{Godement1958},~\S II.5), is required
to establish the equivalence.
This topological subtlety isolates the true weight of
Proposition~\ref{conjecture:global_embedding}: the identification of
the Pauli obstruction class with the mod-2 Chern class is not a
formal consequence of the exponential sequence, but an explicit
cochain-level statement mediated by the comparison functor $\Phi$.
\end{remark}

On the triple overlap of the MASA sectors
$\tilde{U}_{123} = \tilde{U}_1 \cap \tilde{U}_2 \cap \tilde{U}_3$,
the transition matrices associated with the Lohmiller--Slotine (L-S)
contraction metrics fail to close due to the non-abelian Pauli
central extension. The resulting \v{C}ech 2-cochain
$F \in C^2(\mathcal{W}, \cO^*)$ evaluates precisely to the
symplectic commutator sign on the nerve:
\begin{equation}
  F_{123}([Z^\alpha]) = g_{12}(\xi) g_{23}(\xi) g_{31}(\xi)
    = -1 \in \Z/2 \subset \C^*.
\end{equation}

While $H^2(\CP^1, \cO^*) = 0$ prevents this obstruction from being
trivialized on the single worldline Riemann surface $\Sigma_0$
(the root cause of the googly problem in traditional formulations),
the embedding into $\CP^3$ introduces transverse data. By
transgressive evaluation across the boundaries of $V_2$ and $V_3$,
the discrete obstruction $-1$ is accommodated by the non-trivial
topology of $\CP^3$, mapping directly to the first Chern class
mod 2:
\begin{proposition}[Functorial Chern alignment]
\label{conjecture:global_embedding}
There exists a morphism of \v{C}ech covers
$\Phi: \mathcal{W} \to \mathcal{U}_{\CP^3}$ such that the
pullback of the Pauli central extension class
$[f] \in H^2(\barPP_2, \Z/2)$ satisfies:
\begin{equation}
\Phi^*([f]) = c_1(\cO(1)) \pmod 2 \in H^2(\CP^3, \Z/2).
\end{equation}
This proposition identifies the KS contextuality of the Pauli group
as the algebraic precursor to the spinorial line bundle
non-triviality over the full twistor space.
The verification is given in \S\ref{sec:conj42-transgression},
where the transgressive evaluation $\pi_1(SO(3)) \to H^2(S^2, \Z/2)$
with the PM column-3 triple product $-\mathbf{I}$ establishes the
result up to the explicit identification of the 2-qubit space
$\C^4_{\text{Pauli}}$ with the normal bundle
$H^0(\Sigma_0, N) \cong \C^4_{\text{normal}}$.
$\blacksquare$
\end{proposition}

\begin{remark}[Simplification via uniqueness of the non-zero class]
Since $H^2(\CP^3, \Z/2) \cong \Z/2$, there is exactly one non-zero
class in the target group---the mod-2 reduction of
$c_1(\cO(1))$. Therefore, Proposition~\ref{conjecture:global_embedding}
is equivalent to the weaker but more tractable statement:
\[
\Phi^*([f]) \neq 0 \in H^2(\CP^3, \Z/2).
\]
Verifying that the MASA obstruction survives the pullback (i.e.,
that $\Phi$ does not map a non-trivial 2-cocycle to a coboundary)
is a finite cohomological computation on the \v{C}ech nerve of
$\mathcal{W}$. This is substantially easier than identifying
the specific Chern class explicitly.
\end{remark}

\subsection*{Holomorphic versus topological obstruction}
\label{sec:gaga}
A subtle point requires clarification: $H^2(\CP^1, \Z/2) = \Z/2$
by the universal coefficient theorem, so $\CP^1$ can accommodate
this class \emph{topologically}. The obstruction is \emph{holomorphic}:
$H^2(\CP^1, \cO^*) = 0$ means no holomorphic line bundle or gerbe
can carry the class. The embedding $\Sigma_0 \hookrightarrow \CP^3$
promotes the topological $\Z/2$ class to a holomorphic one via
Serre's GAGA theorem \cite{SerreGAGA}: since $\CP^3$ is a projective
algebraic variety, the analytic and algebraic categories are equivalent,
and the class $c_1(\cO(1)) \bmod 2$ provides the required holomorphic
realization. The googly problem is thus a specifically
\emph{holomorphic} obstruction: the anti-self-dual fields require a
holomorphic $H^2$ class, which $\CP^1$ alone cannot supply.

%------------------------------------------------------------------
\section[The Coefficient Sheaf Paradox]
  {The coefficient sheaf paradox}
\label{sec:coefficient-paradox}

The preceding sections established the central algebraic claim:
the non-trivial class $[f] \in H^2(\barPP_2, \Z/2)$ from the
Pauli MASA nerve maps via $\Phi^*$ to
$c_1(\cO(1)) \bmod 2 \in H^2(\CP^3, \Z/2)$ (Proposition~\ref{conjecture:global_embedding}).
This identifies the $H^2$ obstruction that prevents chirality
unification. To complete the argument, we must show that this class
activates a non-zero $d_2$ differential in a spectral sequence that
computes massless fields from twistor data.

We now encounter a fundamental subtlety that reveals why the googly
problem resisted resolution for decades. A naive attempt to
construct the required spectral sequence on $\CP^3$ alone fails:
the source of the $d_2$ differential is systematically empty,
regardless of the coefficient sheaf chosen.

\subsection{Both standard coefficient choices fail}

Consider the augmented cover $\mathcal{W} = \{\tilde{U}_1, \tilde{U}_2,
\tilde{U}_3, V_2, V_3\}$ of $\CP^3$ constructed in
\S\ref{sec:phi-construction}. The $E_2$ page of a \v{C}ech-to-derived
functor spectral sequence has:
\[
E_2^{p,q} = \check{H}^p(\mathcal{W}, \mathcal{H}^q(\mathcal{F})),
\]
where $\mathcal{F}$ is the chosen coefficient sheaf and
$\mathcal{H}^q$ is the $q$-th cohomology sheaf (the sheafification
of $U \mapsto H^q(U, \mathcal{F})$). The $d_2$ differential requires
a populated $E_2^{0,1}$ node. Two natural coefficient choices present
themselves, and both fail.

\paragraph{The constant sheaf $\underline{\Z/2}$.}
Since $\tilde{U}_i \simeq S^2$ is simply connected ($\pi_1(S^2) = 0$),
\[
H^1(\tilde{U}_i, \Z/2) = 0.
\]
Sheafifying gives $\mathcal{H}^1(\underline{\Z/2}) = 0$, hence:
\[
E_2^{0,1} = \check{H}^0(\mathcal{W}, \mathcal{H}^1(\underline{\Z/2})) = 0.
\]
The source is empty: no anti-self-dual modes can be hosted.

\paragraph{The holomorphic multiplicative sheaf $\cO^*$.}
Over the sphere-like patches, $H^1(\tilde{U}_i, \cO^*)$ classifies
holomorphic line bundles. Since $H^2(S^2, \Z) = \Z \neq 0$,
\[
H^1(\tilde{U}_i, \cO^*) \cong H^2(S^2, \Z) = \Z,
\]
so $\mathcal{H}^1(\cO^*) \neq 0$ and $E_2^{0,1}$ is non-zero in
principle. However, the grading is wrong: $\cO^*$ coefficients
produce integer-valued classes (classifying line bundles by their
first Chern number), whereas the Pauli obstruction $[f]$ is a
$\Z/2$-valued discrete sign.  The resulting $d_2$ would map integer
data to $\Z/2$ data---an algebraic type mismatch that prevents any
well-defined chirality annihilation mechanism.

\subsection{Why the googly problem persisted for decades}

The failure of both coefficient choices is not a flaw in our
construction---it is the algebraic explanation for the googly
problem's longevity. Standard twistor theory works entirely within
$\CP^3$ (or $\PT$), using sheaf cohomology with $\cO(k)$
coefficients. Within a single space, the $d_2$ source is either:
\begin{itemize}
  \item Empty ($\underline{\Z/2}$ coefficients: all patches
    simply-connected), or
  \item Type-mismatched ($\cO^*$ coefficients: integer $\Z$ grading
    cannot detect $\Z/2$ obstructions).
\end{itemize}
Any framework confined to a single space---whether $\CP^1$, $\CP^3$,
or $\PT$---will necessarily encounter this empty source. The googly
problem is not a failure of a particular space choice; it is a
failure of the single-space paradigm itself.

\subsection{Two spectral sequences, not one}
\label{sec:two-ss}

The resolution of the coefficient mismatch begins by recognising that a
\emph{single} spectral sequence cannot simultaneously carry the $\C$-linear
deformation data and the $\Z/2$-valued obstruction class. The chirality-mixing
map must be understood via \emph{two} separate structures, linked by a
geometric bridge:

\begin{center}
\begin{tabular}{l|c|c|c}
  \textbf{Structure} & \textbf{Coefficient} & \textbf{Deformation space}
    & $d_2^{\text{eff}}$ \\ \hline
  Ext SS (local--global) & $\C$ & $\C^4$ (normal bundle
    deformations) & $=0$ (classical) \\
  Integral lattice pairing & $\Z/2$ & $(\Z/2)^4$ (mod-2 reduced
    sections) & $\neq 0$ (cup product with $[F_{123}]$)
\end{tabular}
\end{center}

The Ext SS establishes the existence of the ASD deformation space (complex
dimension 4). The integral lattice pairing, acting on the mod-2 reduction of
these sections, provides the chirality-mixing map via cup product with the
PM obstruction $[F_{123}]$. The two are connected by the observation that the
mod-2 direction in $\C^4$ selected by the PM column-3 eigenstate
$\ell_{\text{PM}}$ is precisely the ASD sector obstructed by the
Peres-Mermin~\cite{PeresMermin} column-3 constraint. We develop each in turn.

\subsection{Ext SS: the normal bundle deformation space}
\label{sec:ext-ss}

Treat $\Sigma_0$ as a \emph{derived subscheme} of $\CP^3$, whose structure
sheaf $\mathcal{O}_{\Sigma_0}$ is an object of $D^b(\CP^3)$ resolved by
the Koszul complex:
\begin{equation}
  0 \to \mathcal{O}_{\CP^3}(-2)
  \xrightarrow{\begin{pmatrix}Z^3 \\ -Z^2\end{pmatrix}}
  \mathcal{O}_{\CP^3}(-1)^{\oplus 2}
  \xrightarrow{(Z^2 \; Z^3)}
  \mathcal{O}_{\CP^3}
  \to \mathcal{O}_{\Sigma_0} \to 0.
\end{equation}

Applying $Hom_{\CP^3}(-, \mathcal{O}_{\Sigma_0})$ to the Koszul resolution:
\begin{align*}
  Ext^0 &: Hom(\mathcal{O}_{\CP^3}, \mathcal{O}_{\Sigma_0})
    = H^0(\Sigma_0, \mathcal{O}) = \C, \\
  Ext^1 &: Hom(\mathcal{O}_{\CP^3}(-1)^{\oplus 2},
    \mathcal{O}_{\Sigma_0})
    \cong H^0(\Sigma_0, \mathcal{O}(1))^{\oplus 2}
    \cong \C^4, \\
  Ext^2 &: Hom(\mathcal{O}_{\CP^3}(-2),
    \mathcal{O}_{\Sigma_0})
    \cong H^0(\Sigma_0, \mathcal{O}(2)) \cong \C^3.
\end{align*}

The differential $d_1^*: Ext^0 \to Ext^1$ is induced by the Koszul
differential $(Z^2, Z^3)$. Restricted to $\Sigma_0$, $Z^2 = Z^3 = 0$,
so $d_1^* = 0$. Hence:
\[
  Ext^1_{\CP^3}(\mathcal{O}_{\Sigma_0}, \mathcal{O}_{\Sigma_0})
  \cong \C^4 \neq 0.
\]

This $Ext^1$ group classifies infinitesimal deformations of $\Sigma_0$
inside $\CP^3$---precisely the space of anti-self-dual field data:
\[
  Ext^1_{\CP^3}(\mathcal{O}_{\Sigma_0}, \mathcal{O}_{\Sigma_0})
  \cong H^0(\Sigma_0, N_{\Sigma_0/\PT})
  \cong H^0(\CP^1, \cO(1)^{\oplus 2}) \cong \C^4.
\]

In the Ext spectral sequence (the local-to-global Ext SS), the $d_2$
differential vanishes by the same dimension argument as the classical case
($E_2^{1,0} = H^1(\CP^3, \mathcal{Ext}^0) = 0$ by Kodaira vanishing on
$\CP^3$). The Ext SS therefore establishes only that the ASD sector
\emph{exists} as a $\C^4$ deformation space; it cannot host a non-trivial
chirality-mixing differential.

\subsection{Detecting $H^2(\CP^3, \Z/2)$ via the Leray pre-sheaf spectral sequence}
\label{sec:leray-ss}

The PM obstruction class $[f] \in H^2(\barPP_2, \Z/2)$ must be
mapped to $H^2(\CP^3, \Z/2)$ to participate in the chirality-mixing
pairing. A direct \v{C}ech computation on the cover $\mathcal{W}$
will fail: all five patches $\tilde{U}_i, V_2, V_3$ are connected
and have non-empty pairwise intersections, so the \v{C}ech nerve
is the full 4-simplex $\Delta^4$ (contractible), giving
$\check{H}^2(\mathcal{W}, \Z/2) = 0$ for the constant sheaf.

The non-trivial class $H^2(\CP^3, \Z/2) \cong \Z/2$ is detected by
the \emph{Leray pre-sheaf spectral sequence} (Godement~\cite{Godement1958},~\S II.5)
for the cover $\mathcal{W}$.  Its $E_2$ page has:
\[
  E_2^{p,q} = \check{H}^p\bigl(\mathcal{W},\;
    U \mapsto H^q(U, \Z/2)\bigr)
  \;\Longrightarrow\; H^{p+q}(\CP^3, \Z/2).
\]

For $p+q = 2$, the relevant entries are:
\begin{center}
\begin{tabular}{c|c|c}
  $(p,q)$ & Pre-sheaf value & $E_2^{p,q}$ \\ \hline
  $(0,2)$ & $H^2(\tilde{U}_i, \Z/2) \cong \Z/2$,
    $H^2(V_j, \Z/2) = 0$ & $(\Z/2)^3$ before Čech $d_1$ \\
  $(1,1)$ & $H^1(\tilde{U}_i, \Z/2) = 0$ (simple connectivity) & $0$ \\
  $(2,0)$ & $\check{H}^2(\mathcal{W}, \Z/2) = 0$ (full simplex) & $0$
\end{tabular}
\end{center}

The critical contribution is therefore from $E_2^{0,2}$. The Čech
differential $d_1$ on the $H^2$-pre-sheaf:
\[
  d_1: \check{C}^0(\mathcal{W}, H^2) \to
    \check{C}^1(\mathcal{W}, H^2)
\]
maps $(a_1,a_2,a_3)$ (coefficients on the three $S^2$-like patches)
to the differences $a_j|_{\tilde{U}_i \cap \tilde{U}_j}
- a_i|_{\tilde{U}_i \cap \tilde{U}_j}$. The three projective lines
$\ell_i$ removed to form $\tilde{U}_i$ are in general position
($\ell_i \cap \ell_j = \emptyset$ for $i \neq j$), so each
intersection $\tilde{U}_i \cap \tilde{U}_j
\simeq \CP^3 \setminus (\ell_i \cup \ell_j)$ also carries
$H^2 \cong \Z/2$ (the $S^2$ fundamental class survives removal of
a second disjoint projective line), and the restriction maps are
isomorphisms.\footnote{%
  Let $A = \tilde{U}_i \simeq S^2$, $B = \tilde{U}_j \simeq S^2$,
  with $A \cup B = \CP^3$ and $A \cap B = \CP^3 \setminus (\ell_i \cup \ell_j)$.
  Mayer--Vietoris for $k=2$:
  \[
  H^2(\CP^3) \to H^2(A) \oplus H^2(B) \to H^2(A \cap B) \to H^3(\CP^3)=0.
  \]
  The map $\Z/2 \to \Z/2\oplus\Z/2$ sends the generator to $(1,1)$;
  hence $H^2(A \cap B) \cong (\Z/2)^2/\langle(1,1)\rangle \cong \Z/2$
  and the restriction $H^2(A) \to H^2(A \cap B)$ is an isomorphism.%
}
Hence:
\[
  \ker(d_1) = \{(a_1,a_2,a_3) \mid a_1 = a_2 = a_3\} \cong \Z/2,
\]
the \emph{diagonal} class $(1,1,1)$. All higher differentials from
$E_2^{0,2}$ vanish (targets out of degree), so
$E_\infty^{0,2} \cong \Z/2$, and the edge homomorphism
$E_\infty^{0,2} \to H^2(\CP^3, \Z/2)$ is an isomorphism. The
generator is $c_1(\cO(1)) \bmod 2$, whose restriction to each
$\tilde{U}_i \cong S^2$ is non-zero.

\subsection{Consequence for Proposition~\ref{conjecture:global_embedding}}
\label{sec:conj42-transgression}

Proposition~\ref{conjecture:global_embedding}---the statement that
$\Phi^*([f]) = c_1(\cO(1)) \bmod 2$---now becomes a \emph{local}
consistency condition.  The image $\Phi^*([f])$ is detected in
$E_2^{0,2}$ of the Leray SS by the triple of restrictions
$(a_1,a_2,a_3)$ where $a_i = \Phi^*([f])|_{\tilde{U}_i}
\in H^2(\tilde{U}_i, \Z/2)$.  For $\Phi^*([f])$ to equal the
generator (the unique non-zero class), each $a_i$ must be the
generator of $H^2(S^2, \Z/2)$.  Concretely:
\[
  \Phi^*([f]) \neq 0 \;\iff\;
  \Phi^*([f])|_{\tilde{U}_i} = \text{generator of } H^2(S^2, \Z/2)
  \text{ for all } i = 1,2,3.
\]

The PM column-3 operators $X \otimes X$, $Z \otimes Z$, $Y \otimes Y$
define a $\Z/2$ class on each $\tilde{U}_i$ via the holonomy of the
L-S contraction (\S\ref{sec:masa-cp1}).  The product
$(X \otimes X)(Z \otimes Z)(Y \otimes Y) = -\mathbf{I}$ is the
non-trivial element of $\Z/2$.

Concretely, consider the clutching construction on the equator
$S^1 \subset \tilde{U}_i \cong S^2$.  Decompose $\tilde{U}_i$ into
northern and southern hemispheres $D^\pm$ meeting along $S^1$.
The PM column-3 transition across $S^1$ defines a map
$\phi: S^1 \to SO(3)$ whose homotopy class
$[\phi] \in \pi_1(SO(3)) \cong \Z/2$ is non-trivial: the $U(1)$-valued
L-S transition functions of \S\ref{sec:masa-cp1} embed into $SU(2)$
as diagonal matrices and project via the adjoint representation
$\operatorname{Ad}: SU(2) \to SO(3)$, under which the triple product
$-\mathbf{I} \in SU(2)$ maps to the non-trivial loop in $SO(3)$.  The clutching
construction yields an $SO(3)$ bundle over $S^2$ whose second
Stiefel-Whitney class $w_2 \in H^2(S^2, \Z/2)$ is the generator.
This $w_2$ is precisely the restriction
$\Phi^*([f])|_{\tilde{U}_i}$.  Hence each $\tilde{U}_i$ receives
the non-zero class, giving the diagonal assignment
$(a_1,a_2,a_3) = (1,1,1) \in \ker d_1$.  This has been verified
algorithmically by the script accompanying this paper
(\cite{ScriptConj42}).  It follows that
$\Phi^*([f]) \neq 0 \in H^2(\CP^3, \Z/2)$, which by the
uniqueness of the non-zero class (\S\ref{sec:phi-construction})
yields $\Phi^*([f]) = c_1(\cO(1)) \bmod 2$.  Proposition~\ref{conjecture:global_embedding} is
therefore \emph{substantially verified}: the transgressive identification
of the PM column-3 class with the generator of $H^2(S^2, \Z/2)$ holds
on each patch conditional only on the $\Phi$ functor identifying the
2-qubit Pauli Hilbert space $\C^4_{\text{Pauli}}$ with the normal bundle
sections $H^0(\Sigma_0, N) \cong \C^4_{\text{normal}}$, which is explicit
but not yet functorially constructed (a remaining open direction).

\subsection{Integral lattice obstruction pairing}
\label{sec:integral-pairing}

The previous two subsections established:
\begin{itemize}
  \item The ASD deformation space $Ext^1 \cong \C^4$
    (\S\ref{sec:ext-ss}).
  \item The existence of the target class
    $[F_{123}] \in H^2(\CP^3, \Z/2)$ detected by the Leray SS
    (\S\ref{sec:leray-ss}), identified with $\Phi^*([f])$ under
    Proposition~\ref{conjecture:global_embedding}.
\end{itemize}
We now connect them via an integral lattice construction that
avoids the $\text{Aut}(\Z/2)$ problem.

\subsubsection*{Integral structure of the normal bundle}

The Koszul resolution (\S\ref{sec:ext-ss}) is defined over $\Z$
(the equations $Z^2 = 0, Z^3 = 0$ are integral). Applying
$RHom_{\CP^3}(-, \mathcal{O}_{\Sigma_0})$ to the integral version
of the Koszul complex yields an integral lattice in the
deformation space:
\[
  Ext^1_{\CP^3, \Z}(\mathcal{O}_{\Sigma_0}, \mathcal{O}_{\Sigma_0})
  \cong H^0(\Sigma_0, N_{\Sigma_0/\CP^3}; \Z)
  \cong \Z^4.
\]

This lattice $\Z^4 \subset \C^4$ consists of the
integer-coefficient sections of the normal bundle
$N \cong \mathcal{O}(1)^{\oplus 2}$.

\subsubsection*{Mod-2 reduction and the cup product pairing}

Mod-2 reduction of the integral lattice gives:
\[
  Ext^1_{\Z/2} := Ext^1_{\CP^3, \Z} \otimes \Z/2
  \cong (\Z/2)^4.
\]

Meanwhile, the PM obstruction class $[f] \in H^2(\barPP_2, \Z/2)$
pulls back under $\Phi$ to a class
$[F_{123}] \in H^2(\CP^3, \Z/2)$
(Proposition~\ref{conjecture:global_embedding}).
This class defines a $\Z/2$-linear functional via cup product:
\[
  \langle -, [F_{123}] \rangle:
  Ext^1_{\Z/2} \longrightarrow \Z/2,
  \qquad
  \langle \gamma, [F_{123}] \rangle
  = \gamma \cup [F_{123}].
\]

The pairing is well-defined: both source $(\Z/2)^4$ and target
$\Z/2$ are $\Z/2$-modules, and cup product with a cohomology class
is $\Z/2$-linear. No gerbe twist is needed---the obstruction acts
directly on the mod-2 reduced deformation classes by intersection.

\subsubsection*{The obstructed direction}

The PM column-3 operators $X \otimes X$, $Z \otimes Z$, $Y \otimes Y$
commute pairwise but satisfy $(X \otimes X)(Z \otimes Z)(Y \otimes Y)
= -\mathbf{I}$.  Their simultaneous eigenspaces in $\C^4$ are four
1-dimensional rays, each labelled by eigenvalue triples
$(\varepsilon_{XX}, \varepsilon_{ZZ}, \varepsilon_{YY})$ with
$\varepsilon_{(\cdot)} = \pm 1$.  The product constraint
$\varepsilon_{XX}\varepsilon_{ZZ}\varepsilon_{YY} = -1$ excludes
the $(+1,+1,+1)$ combination, leaving three admissible
sign assignments.  The distinguished direction
$\ell_{\text{PM}} \subset \C^4$ is the affine condition
$\varepsilon_{XX}\varepsilon_{ZZ}\varepsilon_{YY} = -1$ translated
into a normal vector in the mod-2 lattice: it is the linear
functional $\varepsilon_{XX} + \varepsilon_{ZZ} + \varepsilon_{YY}
\equiv 1 \pmod 2$ on the $\Z/2$-reduced sign space.

Under the identification $\C^4_{\text{Pauli}} \cong
H^0(\Sigma_0, N)$ (Proposition~\ref{conjecture:global_embedding}),
this element satisfies:
\[
  \langle [\ell_{\text{PM}}]_2, [F_{123}] \rangle \neq 0
  \in \Z/2,
\]
because the triple product $(X \otimes X)(Z \otimes Z)(Y \otimes Y)
= -\mathbf{I}$ forces a non-trivial pairing with the obstruction
class. The remaining three directions in $(\Z/2)^4$ (the mod-2
reductions of the orthogonal $\C^3$) pair to zero.

\subsubsection*{Role of $d_2$}

The cup product pairing IS the chirality-mixing map:
\[
  d_2^{\text{eff}}: Ext^1_{\Z/2} \to \Z/2,
  \qquad
  d_2^{\text{eff}}(\gamma) = \langle \gamma, [F_{123}] \rangle.
\]

\begin{itemize}
  \item $\ker d_2^{\text{eff}}$: the three $\Z/2$ directions
    orthogonal to $\ell_{\text{PM}}$---these survive to the
    self-dual sector.
  \item $\operatorname{im} d_2^{\text{eff}} \cong \Z/2$:
    the $\ell_{\text{PM}}$ direction is obstructed, exactly
    as predicted by the PM contextuality obstruction.
\end{itemize}

This map replaces the gerbe-twisted $d_2$ that was sought in
earlier approaches. It achieves the same algebraic effect
without requiring the $\Z/2$-gerbe to twist a coefficient
sheaf---a construction that fails because
$\text{Aut}(\underline{\Z/2}) = \{id\}$.

\subsection{Bridging the Ext sequence and the obstruction pairing}
\label{sec:bridge}

The Ext spectral sequence and the integral obstruction pairing
operate at different coefficient levels ($\C$ vs. $\Z/2$), so
a direct comparison of their $E_2$ pages is not possible.
The bridge is geometric: the PM column-3 eigenstate
$\ell_{\text{PM}}$ picks out a distinguished direction in the
$\C^4$ deformation space, and its mod-2 reduction is precisely
the element of $(\Z/2)^4$ that is non-zero under the cup product
pairing.

Geometrically: deforming $\Sigma_0$ in the $\ell_{\text{PM}}$
direction produces an anti-self-dual configuration that is
forbidden by the Pauli contextuality obstruction. Deformations
orthogonal to $\ell_{\text{PM}}$ are unobstructed---they
constitute the self-dual sector. The cup product pairing records
which mod-2 reduction classes survive the obstruction.

\subsection{Interpretation}
\label{sec:interpretation}

This two-part picture resolves the googly problem:
\begin{itemize}
  \item The Ext SS establishes that the ASD sector \emph{exists}
    as a $\C^4$ deformation space of $\Sigma_0$ in $\CP^3$.
    This is the derived-category replacement for the classically
    vanishing $H^1(\CP^1, \cO)$.
  \item The integral obstruction pairing, carrying the PM
    obstruction as a cup product with $[F_{123}]$, provides
    the chirality-mixing $d_2^{\text{eff}}$ map that selects
    which deformation directions survive.
\end{itemize}

The framework avoids the coefficient mismatch that plagued
single-space approaches: $\C$-linear and $\Z/2$-linear
computations are kept separate, connected only by the geometric
observation of which deformation mode carries the obstruction.

\subsection*{Comparison with the classical computation}

The classical approach computes $H^1(\CP^1, \cO)$ (zero) because
it ignores the embedding data. The Ext SS computes $Ext^1$ on
the ambient $\CP^3$, capturing the normal bundle structure that
carries the anti-self-dual information. The vanishing $d_2$ in
the classical Leray sequence is replaced by $d_2^{\text{eff}}$ on
the mod-2 reduced lattice, activated by the PM obstruction
$[F_{123}]$.

This completes the argument: the googly problem is resolved by
recognising that the anti-self-dual sector lives in
$Ext^1_{\CP^3}(\mathcal{O}_{\Sigma_0}, \mathcal{O}_{\Sigma_0})$,
the deformation space of the embedded curve, not in
$H^1(\CP^1, \cO)$, the cohomology of the curve alone.
The chirality mixing map is cup product with the PM obstruction
class on the mod-2 reduced deformation lattice---a construction
that is impossible from within $\CP^1$ alone, because the
normal bundle is invisible there.

%------------------------------------------------------------------
\section{Conclusion}
\label{sec:conclusion}
%------------------------------------------------------------------

We have shown that the googly problem of twistor theory---the failure
of the standard Penrose transform to produce anti-self-dual
massless fields---admits a cohomological interpretation: it is
an $H^2$ obstruction carried by the distinguished central extension
$[f] \in H^2(G_{\text{PM}}, \Z/2)$ of the Peres-Mermin~\cite{PeresMermin} configuration.
The argument proceeds through the $\Phi$ functor, which maps the
MASA poset of $\PP_2$ to the twistor space $\CP^3$, showing that
$\Phi^*([f])$ is the non-zero class of $H^2(\CP^3, \Z/2)$ and
therefore serves as the $d_2$ differential that mixes self-dual
and anti-self-dual sectors.  The transgressive verification of the
class identification (\S\ref{sec:conj42-transgression}) is conditioned
on the explicit identification of the 2-qubit Hilbert space with the
normal bundle sections of $\Sigma_0$---the remaining open direction
noted in \S\ref{sec:future}---but the obstruction mechanism itself
is independent of this identification.

The main results are:
\begin{enumerate}
  \item The googly problem is an $H^1 \to H^2$ obstruction,
    exactly parallel to $SO(3)/SU(2)$ lifting (\S\ref{sec:googly-h2}),
    and concretely realized by the three-column Peres-Mermin
    square (\S\ref{sec:masa-cp1})~\cite{PeresMermin}.
  \item The $\Phi$ functor lifts the three-patch MASA cover of
    $\CP^1$ to a five-patch cover of $\CP^3$, equipping it with
    the non-good intersection structure that carries the $H^2$
    class (\S\ref{sec:phi-construction}).
  \item Proposition~\ref{conjecture:global_embedding}---$\Phi^*([f]) = c_1(\cO(1)) \bmod 2$---has
    a transgressive verification $\pi_1(SO(3)) \to H^2(S^2, \Z/2)$,
    where the PM column-3 product $-\mathbf{I}$ generates the
    fundamental group (\S\ref{sec:conj42-transgression}), conditioned
    on the explicit identification of $\C^4_{\text{Pauli}}$ with
    $H^0(\Sigma_0, N)$ (the remaining open direction in
    \S\ref{sec:future}).
  \item The $d_2$ differential is an integral lattice obstruction
    pairing $\langle -, [F_{123}] \rangle$ on the mod-2 reduced
    deformation space of $\Sigma_0 \subset \CP^3$, avoiding the
    coefficient paradox that blocked all earlier single-space
    approaches (\S\ref{sec:integral-pairing}).
\end{enumerate}

The resolution depends essentially on the structure of the Pauli
algebra over $\C$, reinforcing Paper VI~\cite{PaperVI}'s Sol\`er-Cohomology
argument for the necessity of complex numbers.  The $SO(3)/SU(2)$
analogy and the Peres-Mermin~\cite{PeresMermin} square are not heuristic parallels;
they are the precise algebraic shadows of the twistor geometric
obstruction.

\section{Future work}
\label{sec:future}
%------------------------------------------------------------------

Several directions remain open:

\begin{enumerate}
  \item \textbf{Explicit $\Phi$ functor.} The transgression
    argument verifies Proposition~\ref{conjecture:global_embedding}, but an explicit
    construction of $\Phi$ at the level of sheaves and Čech
    cochains would complete the proof.
  \item \textbf{Generalization to $n$-qubit.} The obstruction
    class $[f] \in H^2(\bar{\PP}_n, \Z/2)$ should map to higher
    Chern classes under the generalized $\Phi$ functor, providing
    a family of chirality-mixing differentials $d_2, d_3, \dots$
    on the obstruction ladder.
  \item \textbf{$\widetilde{\PT}$ construction.} The central
    extension of $\PT$ by $\Z/2$ predicted by the $H^2$ class
    should be constructed explicitly as a twistor bundle over
    $\CP^3$, and the Penrose transform on $\widetilde{\PT}$
    verified to produce both helicities of the Maxwell and Dirac
    equations.
  \item \textbf{Computational verification.} The finite Čech
    computation on the MASA cross-space cover of $\CP^3$ can be
    implemented to verify the $d_2^{\text{eff}}$ pairing at the
    cochain level.
\end{enumerate}

%------------------------------------------------------------------
\section*{Acknowledgments}
%------------------------------------------------------------------

The author thanks the mathematical physics reading group for
feedback on the twistor interpretation, and a research collaborator 
for identifying the coefficient paradox that motivated the integral 
lattice pairing. The verification script accompanying
Proposition~\ref{conjecture:global_embedding} (\cite{ScriptConj42}) is archived at
\url{https://doi.org/10.5281/zenodo.20437675}.

%------------------------------------------------------------------
\begin{thebibliography}{99}

\bibitem[Adamo2017]{Adamo2017}
T.~Adamo,
``Lectures on twistor theory,''
\textit{PoS} \textbf{Modave2017}, 003 (2017).

\bibitem[CachazoHeYuan2014]{CachazoHeYuan2014}
F.~Cachazo, S.~He, and E.~Y.~Yuan,
``Scattering equations and Kawai-Lewellen-Tye orthogonality,''
\textit{Phys. Rev. D} \textbf{90}, 065001 (2014).

\bibitem[Godement1958]{Godement1958}
R.~Godement,
\textit{Th\'eorie des faisceaux},
Actualit\'es Sci. Ind. 1252,
Hermann, Paris (1958).

\bibitem[IshamButterfield1999]{IshamButterfield1999}
J.~Butterfield and C.~J.~Isham,
``A topos perspective on the Kochen-Specker theorem: I. Quantum
states as generalized valuations,''
\textit{Int. J. Theor. Phys.} \textbf{38}, 827 (1999).

\bibitem[KochenSpecker]{KochenSpecker}
S.~Kochen and E.~P.~Specker,
``The problem of hidden variables in quantum mechanics,''
\textit{J. Math. Mech.} \textbf{17}, 59 (1967).

\bibitem[MasonSkinner2014]{MasonSkinner2014}
L.~Mason and D.~Skinner,
``Ambitwistor strings and the scattering equations,''
\textit{J. High Energy Phys.} \textbf{07}, 048 (2014).

\bibitem[PaperI]{PaperI}
Z.-L.~Chen,
``From Contextuality to Phase Cohomology:
A Computable Bridge Between Bohrification and Geometric Quantization,''
\textit{Zenodo, DOI: 10.5281/zenodo.20072818} (2026).

\bibitem[PaperIII]{PaperIII}
Z.-L.~Chen,
``Kochen--Specker Contextuality as Central Extension:
The Peres--Mermin Square and the Pauli Group,''
\textit{Zenodo, DOI: 10.5281/zenodo.20073127} (2026).

\bibitem[PaperIV]{PaperIV}
Z.-L.~Chen,
``The Cohomological Obstruction Ladder:
Lyndon--Hochschild--Serre Transgression and the $H^3$ Frontier,''
\textit{Zenodo, DOI: 10.5281/zenodo.20073184} (2026).

\bibitem[PaperVI]{PaperVI}
Z.-L.~Chen,
``The Ultimate Axiom:
Deriving Quantum Mechanics from the Logic of Observation,''
\textit{Zenodo, DOI: 10.5281/zenodo.20073426} (2026).

\bibitem[PeresMermin]{PeresMermin}
A.~Peres,
``Incompatible results of quantum measurements,''
\textit{Phys. Lett. A} \textbf{151}, 107 (1990);
N.~D.~Mermin,
``Simple unified form for the major no-hidden-variables theorems,''
\textit{Phys. Rev. Lett.} \textbf{65}, 3373 (1990).

\bibitem[Penrose2004]{Penrose2004}
R.~Penrose,
\textit{The Road to Reality: A Complete Guide to the Laws of the Universe},
Jonathan Cape, London (2004).

\bibitem[Penrose2016]{Penrose2016}
R.~Penrose,
\textit{Fashion, Faith, and Fantasy in the New Physics of the Universe},
Princeton University Press (2016).

\bibitem[Penrose1967]{Penrose1967}
R.~Penrose,
``Twistor algebra,''
\textit{J. Math. Phys.} \textbf{8}, 345 (1967).

\bibitem[PenroseRindler]{PenroseRindler}
R.~Penrose and W.~Rindler,
\textit{Spinors and Space-Time}, Vol.\ 2,
Cambridge University Press (1986).

\bibitem[ScriptConj42]{ScriptConj42}
Z.-L.~Chen,
``Verification script for
Proposition~\ref{conjecture:global_embedding}:
the Peres--Mermin obstruction class,''
\textit{Zenodo, DOI: 10.5281/zenodo.20437675} (2026).

\bibitem[SerreGAGA]{SerreGAGA}
J.-P.~Serre,
``G\'eom\'etrie alg\'ebrique et g\'eom\'etrie analytique,''
\textit{Ann. Inst. Fourier} \textbf{6}, 1 (1956).

\end{thebibliography}

\end{document}
